\documentclass[11pt]{article}
\usepackage[margin=1in]{geometry}
\usepackage{amsmath,amssymb,amsthm,mathtools}
\usepackage{graphicx}
\usepackage{microtype}
\usepackage[hidelinks]{hyperref}

\newtheorem{theorem}{Theorem}
\newtheorem{lemma}[theorem]{Lemma}
\newtheorem{remark}[theorem]{Remark}
\newcommand{\T}{\mathbb T}
\newcommand{\N}{\mathbb N}
\newcommand{\eps}{\varepsilon}
\newcommand{\norm}[1]{\left\lVert #1\right\rVert}

\title{A Near-Logarithmic Quantitative Ornstein Lower Bound}
\author{Automated Math Discovery Run \texttt{fa\_banach\_001}, lane 4}
\date{9 August 2026}

\begin{document}
\maketitle

\begin{abstract}
Let $K_N$ be the least constant for which every trigonometric polynomial of
coordinate degree at most $N$ on $\T^2$ satisfies
\[
 \norm{\partial_x\partial_y f}_{1}
 \le K_N\bigl(\norm{\partial_x^2 f}_{1}
                  +\norm{\partial_y^2 f}_{1}\bigr).
\]
Kazaniecki and Wojciechowski proved $K_N\gtrsim\sqrt{\log N}$, recalled the
upper bound $K_N\lesssim\log N$, and conjectured sharp logarithmic growth.  We
prove the near-sharp improvement
\[
                 K_N\gtrsim \frac{\log N}{\log\log N}.
\]
The gain comes from estimating the rational Fourier multipliers blockwise on a
geometrically lacunary Riesz product rather than summing the errors of its
individual Fourier coefficients.
\end{abstract}

\section{Source question and result}

The source's simplest quantitative Ornstein problem concerns the mixed second
derivative versus the two pure second derivatives.  It proves the lower
exponent $1/2$, gives a weak-type argument for the logarithmic upper bound, and
asks whether the latter is sharp; see Figure~\ref{fig:source}.

\begin{figure}[ht]
\centering
\includegraphics[width=0.94\textwidth]{figures/source_question_crop.png}
\caption{The sharpness question in the source paper, arXiv:2206.13666.}
\label{fig:source}
\end{figure}

For definiteness, the degree of a trigonometric polynomial on $\T^2$ is the
maximum absolute value of either coordinate in its Fourier support.  Define
\[
 K_N=\sup_{\substack{f\ne0\\ \deg f\le N}}
 \frac{\norm{\partial_x\partial_yf}_{L^1(\T^2)}}
 {\norm{\partial_x^2f}_{L^1(\T^2)}+
  \norm{\partial_y^2f}_{L^1(\T^2)}}.
\]

\begin{theorem}\label{thm:main}
There are universal constants $c>0$ and $N_0$ such that, for $N\ge N_0$,
\[
                     K_N\ge c\frac{\log N}{\log\log N}.
\]
\end{theorem}

This does not prove the conjectured $K_N\asymp\log N$, but it leaves only one
iterated logarithm between the new lower bound and the known upper bound.

\section{Idea of the proof}

Take lacunary frequencies
\[
        a_k=(M^k,(-1)^kM^k),\qquad 1\le k\le n,
\]
and the Riesz product $P_n=\prod_{k=1}^n(1+\cos\langle a_k,z\rangle)$.
Write $R_n=P_n-1$ as the sum of its martingale-difference blocks $B_k$.
We construct $W_n$ by inverting $\partial_x^2$ on the Fourier support of
$R_n$.  On the $k$th block, the multipliers of the other two derivatives are
\[
 \frac{q_y}{q_x}\approx(-1)^k,
 \qquad
 \left(\frac{q_y}{q_x}\right)^2\approx1.
\]
Thus $\partial_y^2W_n$ is close to $R_n$, whose $L^1$ norm is at most two,
whereas $\partial_x\partial_yW_n$ is close to an alternating transform of the
Riesz blocks, whose $L^1$ norm is comparable to $n$.

The point is the error estimate.  Earlier frequencies contribute a relative
amount $O(1/M)$ to a block frequency.  Expanding the rational multipliers as
power series and using the $L^1$ Bernstein inequality bounds a whole block's
error by $O(1/M)$, without paying for its exponentially many Fourier
coefficients.  Taking $M=n$ makes the total error $O(1)$.  The degree is then
at most $2n^n$, so $n$ is of order $\log N/\log\log N$.

\section{The block multiplier estimate}

Put $\epsilon_k=(-1)^k$, $a_k=(M^k,\epsilon_kM^k)$, and
\[
 P_0=1,\qquad P_k=\prod_{j=1}^k(1+\cos\langle a_j,z\rangle),
 \qquad B_k=P_k-P_{k-1}.
\]
Then
\[
 B_k=\cos\langle a_k,z\rangle P_{k-1},
 \qquad R_n=P_n-1=\sum_{k=1}^nB_k.
\]
Since $M\ge3$, the first coordinate of every nonzero frequency in these
blocks is nonzero.  On their union define Fourier multipliers
\[
 \widehat{T_1h}(q)=\frac{q_y}{q_x}\widehat h(q),
 \qquad
 \widehat{T_2h}(q)=\left(\frac{q_y}{q_x}\right)^2\widehat h(q).
\]

\begin{lemma}[Blockwise error]\label{lem:block}
There is a universal $C$ such that, whenever $M\ge4$,
\[
 \norm{T_1R_n-\sum_{k=1}^n\epsilon_kB_k}_1
 \le C\frac{n}{M-1},
 \qquad
 \norm{T_2R_n-R_n}_1
 \le C\frac{n}{M-1}.
\]
\end{lemma}

\begin{proof}
Fix $k$ and one of the two signed halves of $B_k$.  Its frequencies have the
form
\[
 q=\sigma a_k+s,\qquad \sigma\in\{-1,1\},
 \qquad s\in\operatorname{supp}\widehat P_{k-1}.
\]
Set
\[
 u=\frac{s_x}{\sigma M^k},\qquad
 v=\frac{s_y}{\sigma M^k},\qquad
 \delta_k=\frac{\sum_{j<k}M^j}{M^k}<\frac1{M-1}.
\]
On this half-block,
\[
                 \frac{q_y}{q_x}=\frac{\epsilon_k+v}{1+u}.
\]

We interpret a monomial $u^rv^s$ as a normalized differential multiplier on
$P_{k-1}$.  The one-variable Bernstein inequality in $L^1$, applied
successively in $x$ and $y$, gives
\[
 \norm{u^rv^sP_{k-1}}_1
 \le \delta_k^{r+s}\norm{P_{k-1}}_1
 =\delta_k^{r+s}.
\]
Here $P_{k-1}\ge0$ and its integral is one; lacunarity ensures that its
constant Fourier coefficient is one.

Now expand
\[
 \frac{\epsilon_k+v}{1+u}-\epsilon_k
   =\frac{v-\epsilon_ku}{1+u}
   =\sum_{r\ge0}(-u)^r(v-\epsilon_ku).
\]
The preceding Bernstein estimate and absolute summation show that the
$T_1$ error on the entire $k$th block is at most
\[
                   \frac{2\delta_k}{1-\delta_k}.
\]
For the square multiplier, use
\[
 \left(\frac{\epsilon_k+v}{1+u}\right)^2-1
 =\frac{2\epsilon_kv+v^2-2u-u^2}{(1+u)^2},
 \qquad
 (1+u)^{-2}=\sum_{r\ge0}(-1)^r(r+1)u^r.
\]
Its block error is at most
\[
                  \frac{4\delta_k+2\delta_k^2}{(1-\delta_k)^2}.
\]
The factors $1/2$ from the two signed halves add to one.  Summing these
uniform $O((M-1)^{-1})$ bounds over $k$ proves both assertions.
\end{proof}

\section{The alternating Riesz transform}

We use the same two standard ingredients invoked in the source paper.  To make
the transference precise, write a polynomial on the Riesz spectrum as
\[
 F(x,y)=\sum_{\xi\in\{-1,0,1\}^n}c_\xi
 \exp\!\left(i\sum_{k=1}^n\xi_k(M^kx+\epsilon_kM^ky)\right).
\]
For each fixed $y$, Meyer's lacunary transference theorem applied to the
$x$-frequencies $M^k$ compares $\|F(\cdot,y)\|_{L^1_x}$ with
\[
 \left\|\sum_\xi c_\xi
 \exp\!\left(i\sum_{k=1}^n\xi_k(t_k+\epsilon_kM^ky)\right)
 \right\|_{L^1(\mathbb T^n_t)}.
\]
The latter is independent of $y$, by the Haar translations
$t_k\mapsto t_k-\epsilon_kM^ky$.  Integrating in $y$ therefore compares the
$L^1(\mathbb T^2)$ norm with the fully independent model.  The comparison
constant depends only on the sum of consecutive frequency ratios, and here
\[
 \sum_{k<n}\frac{M^k}{M^{k+1}}=\frac{n-1}{M}.
\]
Hence the comparison constant is universal when $M\ge n$.

In the independent model, the blocks are
\[
 \widetilde B_k=\widetilde P_k-\widetilde P_{k-1},
 \qquad
 \widetilde P_k=\prod_{j=1}^k(1+\cos t_j).
\]
Summation by parts gives
\[
 \sum_{k=1}^n\epsilon_k\widetilde B_k
 =\epsilon_n\widetilde P_n-\epsilon_1\widetilde P_0
  +\sum_{k=1}^{n-1}(\epsilon_k-\epsilon_{k+1})\widetilde P_k.
\]
The lower inequality for linear combinations of finite Riesz products due to
Latała (in the form used by Kazaniecki--Wojciechowski) controls the $L^1$ norm
from below by a universal constant times the total variation of these
coefficients.  Since $|\epsilon_k-\epsilon_{k+1}|=2$, transference yields:

\begin{lemma}[Alternating main term]\label{lem:mainterm}
There is a universal $c_0>0$ such that, for $M\ge n\ge2$,
\[
                       \norm{\sum_{k=1}^n\epsilon_kB_k}_1\ge c_0n.
\]
\end{lemma}

\section{Proof of Theorem~\ref{thm:main}}

Fix $n\ge4$ and set $M=n$.  Define a real trigonometric polynomial $W_n$ by
\[
 \widehat W_n(q)=-\frac{\widehat R_n(q)}{q_x^2}
 \quad(q\in\operatorname{supp}\widehat R_n),
 \qquad \widehat W_n(q)=0\quad\text{otherwise}.
\]
The domination of the last lacunary digit gives
$|q_x|\ge M^k-\sum_{j<k}M^j>0$ on the $k$th block, so this is well defined.
Fourier differentiation gives the exact identities
\[
 \partial_x^2W_n=R_n,
 \qquad \partial_x\partial_yW_n=T_1R_n,
 \qquad \partial_y^2W_n=T_2R_n.
\]
Because $P_n\ge0$ and $\int P_n=1$,
\[
                         \norm{R_n}_1=\norm{P_n-1}_1\le2.
\]
Lemmas~\ref{lem:block} and~\ref{lem:mainterm}, with $M=n$, now imply
\[
 \norm{\partial_y^2W_n}_1\le C_1,
 \qquad
 \norm{\partial_x\partial_yW_n}_1\ge c_0n-C_2.
\]
After increasing the lower threshold for $n$,
\[
 \frac{\norm{\partial_x\partial_yW_n}_1}
 {\norm{\partial_x^2W_n}_1+\norm{\partial_y^2W_n}_1}
 \ge c_1n.
\]

Every Fourier coordinate of $R_n$, and hence of $W_n$, has magnitude at most
\[
                 \sum_{k=1}^nM^k\le2M^n=2n^n.
\]
Thus $K_{2n^n}\ge c_1n$.  Finally, for sufficiently large $N$, take
\[
                  n=\left\lfloor\frac{\log N}{2\log\log N}\right\rfloor.
\]
Then $2n^n\le N$.  Since $K_N$ is nondecreasing,
\[
 K_N\ge c_1n\ge c\frac{\log N}{\log\log N},
\]
which proves the theorem. \qed

\section{Scope and remaining gap}

The proof is specific to the basic isotropic pair
$(\partial_x^2,\partial_y^2)$ and target $\partial_x\partial_y$.  The same
blockwise principle may improve other configurations in the source paper, but
no such generalization is claimed here.  Achieving the conjectured full
$\log N$ lower bound would require uniformly controlling the pure-derivative
error at fixed geometric lacunarity; the present absolute block summation
costs $n/M$ and forces $M\asymp n$, exactly producing the $\log\log N$ loss.

\begin{thebibliography}{9}
\bibitem{KW2024}
K. Kazaniecki and M. Wojciechowski,
\emph{On Bernstein type quantitative estimates for Ornstein
non-inequalities}, Rev. Mat. Iberoam. 40 (2024), 901--912;
arXiv:2206.13666.

\bibitem{Latala2014}
R. Latała,
\emph{$L_1$-norm of combinations of products of independent random
variables}, Israel J. Math. 203 (2014), 295--308.

\bibitem{Meyer1968}
Y. Meyer,
\emph{Endomorphismes des idéaux fermés de $L^1(G)$, classes de Hardy et
séries de Fourier lacunaires}, Ann. Sci. École Norm. Sup. (4) 1 (1968),
499--580.

\bibitem{Bonami2020}
A. Bonami, R. Latała, P. Nayar, and T. Tkocz,
\emph{Bounds on moments of weighted sums of finite Riesz products},
J. Fourier Anal. Appl. 26 (2020), article 17.
\end{thebibliography}

\end{document}
